\documentclass[10pt]{article}
\usepackage[margin=1in]{geometry}
\usepackage[T1]{fontenc}
\usepackage{lmodern,microtype,amsmath,amssymb,amsthm}
\theoremstyle{definition}
\newtheorem{assumption}{Assumption}
\begin{document}
\section*{High-Dimensional Stochastic Nonconvex Stationarity}
We seek an approximate stationary point of a possibly nonconvex function $f:\mathbb{R}^d\to\mathbb{R}$ in a high-dimensional regime. The function is three times continuously differentiable, with globally Lipschitz derivatives of orders one through three. Its value at a given starting point $x^0\in\mathbb{R}^d$ lies within a prescribed gap of its infimum. Each oracle call uses a fresh independent sample to return an unbiased gradient estimate with uniformly bounded mean squared error. For a tolerance $\varepsilon>0$, the goal is a measurable output $\bar{x}\in\mathbb{R}^d$ satisfying $\mathbb{E}\|\nabla f(\bar{x})\|\leq\varepsilon$, averaging over oracle and algorithm randomness.

The underlying optimization problem is
\[
\inf_{x\in\mathbb{R}^d} f(x).
\]
The algorithm receives the dimension $d\in\mathbb{N}$, the starting point $x^0$, the smoothness constants $L_1,L_2,L_3>0$, the noise bound $\sigma\geq 0$, the initial gap bound $\Delta$, and the tolerance $\varepsilon$. To specify the parameter regime, set
\[
M=\min\left\{\frac{L_1}{\varepsilon},\sqrt{\frac{L_2}{\varepsilon}},\sqrt[3]{\frac{L_3}{\varepsilon}}\right\}.
\]
The gap bound and dimension satisfy
\[
\begin{aligned}
\Delta M &\geq 2^{70}\varepsilon,\\
d &\geq 2^{140}\left(1+\frac{\Delta M}{\varepsilon}\left(1+\frac{\sigma^2}{\varepsilon^2}\right)\right)^4.
\end{aligned}
\]

Our assumptions:
\begin{assumption}[Lipschitz derivatives]
The function $f$ belongs to $C^3(\mathbb{R}^d)$. For all $x,y\in\mathbb{R}^d$,
\[
\begin{aligned}
\|\nabla f(x)-\nabla f(y)\| &\leq L_1\|x-y\|,\\
\|\nabla^2 f(x)-\nabla^2 f(y)\|_{\mathrm{op}} &\leq L_2\|x-y\|,\\
\|\nabla^3 f(x)-\nabla^3 f(y)\|_{\mathrm{op}} &\leq L_3\|x-y\|.
\end{aligned}
\]
Here $\|\cdot\|$ is the Euclidean norm and $\|\cdot\|_{\mathrm{op}}$ denotes the induced operator norm.
\end{assumption}

\begin{assumption}[Initial objective gap]
The starting point satisfies
\[
f(x^0)-\inf_{x\in\mathbb{R}^d}f(x)\leq\Delta.
\]
\end{assumption}

\begin{assumption}[Stochastic gradient oracle]
At query $x\in\mathbb{R}^d$ on call $k$, the oracle returns $g(x,\xi_k)$, where $g:\mathbb{R}^d\times\mathbb{R}\to\mathbb{R}^d$ is jointly measurable. The real-valued random variables $(\xi_k)_{k\geq 1}$ are independent with common law $\mu$, and the sequence is independent of the algorithm's randomness. For every $x\in\mathbb{R}^d$ and $\xi\sim\mu$, the vector $g(x,\xi)$ is Bochner integrable, the squared error $\|g(x,\xi)-\nabla f(x)\|^2$ is integrable, and
\[
\begin{aligned}
\mathbb{E}_{\xi\sim\mu}[g(x,\xi)] &= \nabla f(x),\\
\mathbb{E}_{\xi\sim\mu}\bigl[\|g(x,\xi)-\nabla f(x)\|^2\bigr] &\leq \sigma^2.
\end{aligned}
\]
\end{assumption}

After the allotted oracle calls, the output $\bar{x}$ must be measurable and $\|\nabla f(\bar{x})\|$ must be integrable under the joint law $P$ of the algorithm's randomness and the oracle samples. The accuracy requirement is
\[
\mathbb{E}_{P}\bigl[\|\nabla f(\bar{x})\|\bigr]\leq\varepsilon.
\]
\end{document}
